\section{Introduction}
%Drones are beginning to appear in more shared, low-altitude settings rather than only in isolated flight zones. They are increasingly used or proposed for inspection, assistance, education, performance, transportation, parcel delivery, hazardous response, and other tasks near people. In these settings, a drone is not only a remote sensor or flying camera. Its sensing and decisions are tied to a moving physical vehicle that can change the space around it. When drone behavior may affect nearby people, the drone itself, or the surrounding environment, communication becomes part of the safety problem rather than only a matter of convenience. Nearby people therefore need a quick way to express simple or urgent needs, such as asking a drone to move away, hold position, stop a task, or respond to a safety-critical situation, without first finding a controller, opening a mobile phone application, or relying on a networked interface. Spoken input is a natural candidate because voice is immediate, hands-free, and widely treated as a natural modality for human-robot communication~\cite{moore2022vocal}.
Unmanned Aerial Vehicles (UAVs, or drones) are beginning to appear in shared, low-altitude spaces such as classrooms, laboratories, warehouses, inspection sites, and field-support settings. In these environments, a drone may operate autonomously or semi-autonomously while people remain close enough to observe its behavior. Onboard perception, navigation, or task logic may not catch every local risk in time. A nearby person may notice that a drone is drifting toward a wall, approaching a person, or continuing a task that should stop before the drone system resolves the situation on its own. In such moments, human input can become part of the safety picture.

We view the drone in this setting as an embodied computing system: it senses, computes, and acts through a moving physical platform in a space shared with people. A speech safety net therefore supports a boundary for embodied human-drone interaction, where nearby human input may matter before physical consequences occur.

This article proposes speech as a safety net for drone operation. If a drone is about to hit a wall, move too close to a person, or continue a task after the scene has changed, a nearby person should be able to shout ``stop,'' ``move away,'' or ``hold position'' without first finding a controller, opening a mobile phone application, or relying on a networked interface. Speech is a natural candidate for this role because voice is immediate, hands-free, and widely treated as a natural modality for human-robot communication~\cite{moore2022vocal}. The goal is not to replace onboard autonomy or operator control, but to give people in the same space a fast way to express safety-relevant input.

This matters because drone operation has physical consequences. A drone that collides with a wall can damage the vehicle, nearby objects, or the environment; a drone that collides with a person can create injury risk. Regulatory discussions already treat drone operation as a risk-managed activity. Commission Implementing Regulation (EU) 2019/947, for example, frames unmanned aircraft operation around rules and procedures that depend on operating conditions and the protection of people on the ground and other airspace users~\cite{eu2019uasops}. This article argues that, in shared spaces, speech should be considered part of the operating context for safer drone interaction.

We make this idea concrete through drone-side speech processing. Short audio windows are mapped to safety-relevant intent categories: \emph{emergency}, \emph{movement}, and \emph{unknown}. The \emph{emergency} category captures urgent speech that may indicate danger, interruption, or a need to stop. The \emph{movement} category captures ordinary motion-related speech, and \emph{unknown} captures speech or sound that should not be treated as either case.

The article makes three contributions:
\begin{itemize}
    \item It defines speech as a safety-net problem for drone operation in shared low-altitude spaces, where nearby people may observe risks before the drone system resolves them on its own.
    \item It designs a drone-side speech-processing pipeline that maps short audio windows into emergency-related, movement-related, and unsupported categories.
    \item It implements the pipeline on a XIAO ESP32-S3 Sense board and shows that the embedded path can support real-time local audio capture and inference.
\end{itemize}

%However, speech near a drone is not the same as speech to a phone, smart speaker, or cloud assistant. The drone listens from within a changing physical environment: its rotors add noise, its motion changes the distance and angle to the speaker, and nearby people or wind may overlap with the intended utterance. As a result, a phrase may be missed, a background conversation may be treated as relevant, or an urgent stop request may be confused with ordinary movement-related speech. In this setting, a speech error is not only an interface failure. If the error is passed downstream without caution, it may contribute to unsafe behavior, such as a drone continuing a task when someone is trying to interrupt it, moving when it should hold position, or ignoring speech that signals danger to the drone, its surroundings, or nearby people.

%This article argues that spoken interaction with drones needs a safety net between nearby speech and vehicle-facing behavior. The core pipeline is not speech directly becoming action. Instead, nearby audio should first be interpreted by a speech-processing model that can recognize whether the speech is emergency-related, movement-related, or neither. A transcript may be useful in some interfaces, but it is not the required first representation. For drones, the more important step is to turn noisy nearby speech into a compact signal that separates urgent, motion-related, and unsupported speech before it matters to the vehicle.

%We make this safety-net framing concrete through drone-side speech processing. Short audio windows are processed locally and mapped to three categories: \emph{emergency}, \emph{movement}, and \emph{unknown}. The \emph{emergency} category covers urgent speech that may indicate danger, interruption, or a need to stop. The \emph{movement} category covers speech that appears related to ordinary drone motion, such as moving, holding position, or changing direction. The \emph{unknown} category covers speech or sound that should not be treated as either emergency-related or movement-related input. This small vocabulary gives the safety net a practical starting point: urgent speech, ordinary movement-related speech, and input that should not be acted on are kept separate from the beginning.
